\subsection{Сила инерции}
%%%%%%%%%%%%%%%%%%%%%%%%%%%%%%%%%%%%%%%%%%%%%%%%%%%%
%%%%%%%%%%%%%%%%%%%%%%%%%%%%%%%%%%%%%%%%%%%%%%%%%%%% 
 \z[{\begin{wrapfigure}[5]{r}{0.3\textwidth}%
		\vspace{-15pt}%
		\begin{flushright}%
 			\includegraphics[width=0.3\textwidth]{5.5.1}%
		\end{flushright}%
	\end{wrapfigure}%
	}]%
	{Какую постоянную горизонтальную силу нужно приложить к тележке массы $M = 1.0$ кг (см. рисунок), чтобы грузы с массами $m_1 = 0.40$ кг и $m_2 = 0.20$ кг относительно нее не двигались? Трением пренебречь.}
 %
{$F = \frac{(m_1 + m_2 + M) m_2 g}{m_1} = 7.8\,\text{Н}$}
%%%%%%%%%%%%%%%%%%%%%%%%%%%%%%%%%%%%%%%%%%%%%%%%%%%%	
 \z[{\begin{wrapfigure}[4]{r}{0.3\textwidth}%
		\vspace{-15pt}%
		\begin{flushright}%
 			\includegraphics[width=0.3\textwidth]{5.5.2}%
		\end{flushright}%
	\end{wrapfigure}%
	}]%
	{На гладкой горизонтальной поверхности находится призма массы $M$ с углом $\alpha$ при вершине, а на ней брусок массы $m$ (см. рисунок). Коэффициент трения между призмой и бруском равен $\mu$, $\mu > \tan\alpha$. В момент $t = 0$ на призму начала действовать горизонтальная сила, зависящая от времени по закону $F(t) = \beta t$, где $\beta$ — положительная постоянная. Найдите момент времени $t = t_0$, когда брусок начнет скользить по призме.}
 %
{$t_0 = \frac{(m + M)g(\mu - \tan \alpha)}{\beta(1 + \mu \tan \alpha)}$}
%%%%%%%%%%%%%%%%%%%%%%%%%%%%%%%%%%%%%%%%%%%%%%%%%%%%	
 \z[{\begin{wrapfigure}[5]{r}{0.3\textwidth}%
		\vspace{-12pt}%
		\begin{flushright}%
 			\includegraphics[width=0.3\textwidth]{5.5.3}%
		\end{flushright}%
	\end{wrapfigure}%
	}]%
	{На гладкой горизонтальной плоскости находится клин с углом $\alpha$ при основании. Тело массы $m$, положенное на клин, опускается с ускорением, направленным под углом $\beta > \alpha$ к горизонтали. Определите массу клина.}
 %
{$M = m \frac{\tan \alpha}{\tan \beta - \tan \alpha}$}
%%%%%%%%%%%%%%%%%%%%%%%%%%%%%%%%%%%%%%%%%%%%%%%%%%%%
 \z{Горизонтальный диск начинают раскручивать вокруг его оси с линейно возрастающей во времени угловой скоростью $\omega = \varepsilon t$. При какой угловой скорости тело, расположенное на расстоянии $r$ от оси диска, начнет соскальзывать с него, если коэффициент трения между ними равен $\mu$?}
%
{$\omega_1 = \begin{cases}
        0 & \text{при } \varepsilon \geq \frac{\mu g}{R}, \\
        \left(\frac{\mu^2 g^2}{R^2} - \varepsilon^2 \right)^{1/4} & \text{при } \varepsilon < \frac{\mu g}{R}
        \end{cases}$
}
%%%%%%%%%%%%%%%%%%%%%%%%%%%%%%%%%%%%%%%%%%%%%%%%%%%%	
 \z[{\begin{wrapfigure}[6]{r}{0.3\textwidth}%
		\vspace{-10pt}%
		\begin{flushright}%
 			\includegraphics[width=0.3\textwidth]{5.5.5}%
		\end{flushright}%
	\end{wrapfigure}%
	}]%
	{К тяжелому шарику, подвешенному на нити длины $l$, подвешен второй тяжелый шарик на нити той же длины. При вращении шариков вокруг вертикальной оси, проходящей через верхнюю точку подвеса, обе нити лежат в одной плоскости и составляют с вертикалью постоянные углы $\alpha$ и $\beta$. Найдите угловую скорость вращения шариков.}
 %
{$\omega = \sqrt{\frac{g \tan \beta}{l(\sin\beta + \sin\alpha)}}$}